\inicializarSeccion{Conclusiones Individuales}
\inicializarSubseccion{Conclusión de Jared}
\begin{quote}
Esta etapa me ayudó a entender mejor el contexto de nuestro proyecto. Me sorprendió que en Nuevo León dos de cada tres empleos nuevos sean de manufactura. Eso me hizo ver que hay muchísimos técnicos que trabajan con herramientas todos los días y que una solución como ToolNest sí puede servirles. Al analizar Milwaukee TICK, ON!Track y Snap-on, mire que los rastreadores son baratos pero muy limitados, y que las cajas inteligentes funcionan bien pero son muy caras. Creo que nuestra caja puede quedar en un punto medio.
\end{quote}

\inicializarSubseccion{Conclusión de Andrés}
\begin{quote}
Mas que nada, la etapa me ayudo a sentar lo que viene siendo el proyecto en sí. La idea de ToolNest nació precisamente por la necesidad de mantener un orden en mis proyectos. Sin embargo, ví realmente que dicho proyecto tiene el potencial de ayudar a muchisima gente en sectores bastante generales. Sea que quieran mantener un orden personal de sus herramientas o que necesitan matener organizado un taller, la idea de ToolNest es lo suficientemente flexible para poder adaptarse a todos estos escenarios. Espero que en las siguientes etapas estas expectativas se logren cumplir y que mas que un proyecto de clase, que pueda ser utilizado en la vida real.
\end{quote}

\inicializarSubseccion{Conclusión de Elena}
\begin{quote}
Esta etapa me hizo entender las problemáticas que surgen en el mundo del manejo de objetos. Durante esta etapa, me percaté acerca de algo que hasta podría llamar cruel; no es que hayan recursos, es que no sabemos a donde se están yendo. La idea de hacer un prototipo que mantenga la organización de cosas y que además, permita a las personas tener un control realmente inteligente de dichas cosas hará posible que gente cotidiana se avienten a prestar recursos para que se pueda hacer un trabajo más efectivo y de mayor calidad.
\end{quote}

\inicializarSubseccion{Conclusión de Daniel}
\begin{quote}
Partir de los Objetivos de Desarrollo Sostenible (ODS) para identificar una problemática específica y diseñar una solución atractiva y con capacidad real de resolución es un proceso que disfruto. Investigar la problemática en distintos contextos nos permitió comprender que existe una oportunidad genuina en este mercado, y queremos aprovecharla. Tras analizar las soluciones que ya existen, considero que nuestra propuesta es sólida y se diferencia por su enfoque. Además, la integración de distintos sensores y actuadores hará de este prototipo una herramienta muy útil. Estoy seguro de que su desarrollo será un reto enriquecedor y una gran aventura.
\end{quote}
\end{multicols*}